\section{Research Landscape and Literature Gap}

\subsection{Embodied Retrieval and Memory-Augmented Agents}

Embodied retrieval is no longer novel at the conceptual level.
\citet{liu2024embodiedrag} introduce \textit{Embodied-RAG} as a non-parametric memory
that lets embodied agents retrieve past experience to solve new situations.
\citet{li2025embodiedrag} push this direction toward \textit{scene-graph retrieval}
for spatial reasoning and planning in 3D environments.
\citet{xu2024raea} go further by using retrieval over a \textit{policy memory bank}
to support embodied action generation.

The lesson from this line of work is clear: retrieval is already a real
\textit{memory/planning} mechanism in embodied AI, but most existing work focuses on
high-level decision making or action generation rather than reward modeling for continuous manipulation.

\subsection{Retrieval from Demonstrations and Human Video}

Works such as \textit{DINOBot} \citep{johns2024dinobot}, \textit{STRAP} \citep{huang2025strap},
and \textit{Vid2Robot} \citep{jain2024vid2robot} show that retrieval or conditioning with
demonstrations and human videos is a practical direction for robot learning.
\textit{STRAP} is especially relevant here because it suggests that
\textit{sub-trajectory} retrieval can outperform full-episode retrieval for transfer and adaptation.

However, these methods mainly use retrieved evidence to support
\textit{policy execution}, \textit{alignment}, or \textit{behavior cloning}.
Reward is typically secondary rather than the primary contribution.

\subsection{Reward Modeling for Robotics}

Reward modeling in robotics is evolving quickly.
\citet{ma2024vlmfeedback} show that VLM feedback can support
\textit{real-world offline RL} without heavily engineered manual reward.
\citet{yang2025rgvlm} extend this line with \textit{retrieval-guided VLM reward generation}
for robotic offline RL.
By 2026, \citet{song2026roboreward} turn reward modeling itself into a benchmarked research problem with a dedicated robotics reward model.

This line of work proves that reward modeling is meaningful,
but simply saying ``use a multimodal model to synthesize dense reward'' is not novel enough anymore.
The open question is how to connect a reward model with an \textit{external memory}
that is rich enough, stage-aware enough, and reliable enough under OOD conditions.

\subsection{Benchmarks and Simulation Infrastructure}

A simulation-only project is academically valid if it is grounded in the right benchmarks.
\citet{liu2023libero} introduce \textit{LIBERO} for lifelong robot learning and transfer across long-horizon tasks.
\citet{gu2022maniskill2} and \citet{tao2024maniskill3} provide scalable manipulation simulators and efficient data pipelines.
Therefore, a simulation-first scope does not weaken the methodologically central contribution,
as long as the benchmark choice, baselines, and ablations are rigorous.

\subsection{Compact Comparison Table}

\begin{table}[htbp]
  \centering
  \small
  \begin{tabularx}{\textwidth}{p{2.4cm}p{2.7cm}Y Y Y}
    \toprule
    Work family & Representative papers & Is retrieval central? & Is reward modeling central? & Remaining gap relative to the proposed topic \\
    \midrule
    Embodied memory/planning & Embodied-RAG, EmbodiedRAG, RAEA & Yes & Usually no & Strong on planning and memory, but not centered on dense reward for continuous manipulation \\
    Demo/video retrieval & DINOBot, STRAP, Vid2Robot & Yes & Usually no & Strong on action and policy transfer, but does not standardize retrieval-conditioned reward \\
    Robotics reward modeling & VLM Feedback, RG-VLM, RoboReward & Sometimes & Yes & Does not yet fully exploit multi-source external memory for OOD stage-aware reward \\
    Proposed topic & Retrieval-Augmented Reward Modeling & Yes & Yes & Connects multimodal retrieval with reward modeling for long-horizon manipulation in simulation \\
    \bottomrule
  \end{tabularx}
  \caption{Position of the proposed topic relative to the current research landscape.}
\end{table}

\section{Novelty Assessment}

\subsection{What Is Not Novel}

\begin{itemize}
  \item ``Embodied RAG for robots'' is no longer novel as a slogan because retrieval for embodied agents is already an established direction \citep{liu2024embodiedrag, li2025embodiedrag, xu2024raea}.
  \item ``Multimodal reward modeling for robots'' is also no longer an untouched space, since several recent works already use VLMs or LVLMs for robotic feedback and reward \citep{ma2024vlmfeedback, yang2025rgvlm, song2026roboreward}.
\end{itemize}

\subsection{What Can Still Be Novel If the Problem Is Framed Correctly}

\begin{itemize}
  \item Make \textbf{reward modeling} the central contribution, with retrieval acting as a controlled external memory.
  \item Retrieve from \textbf{multiple sources}, including instruction/manual text, stage descriptions, and demonstration clips rather than a single context source.
  \item Emphasize \textbf{stage-aware dense reward} for \textbf{long-horizon} manipulation and evaluate explicitly on \textbf{OOD splits}.
  \item Treat retrieval granularity at the \textbf{sub-trajectory/stage} level as a genuine research variable rather than an implementation detail.
\end{itemize}

\subsection{Novelty Conclusion}

If the topic is framed too broadly, the novelty is only moderate.
If it is framed as:
\begin{quote}
  \textit{Retrieval-Augmented Reward Modeling for Long-Horizon Robot Manipulation in Simulation}
\end{quote}
then the novelty becomes more defensible because the contribution is anchored to
a concrete mechanism, a concrete setting, and a concrete evaluation protocol.

\section{Current Development Directions in the Community}

The current community landscape can be summarized in four directions:
\begin{enumerate}
  \item \textbf{Embodied memory and planning}: retrieval supports reasoning, memory, scene understanding, or planning.
  \item \textbf{Demonstration-conditioned control}: retrieved demonstrations, sub-trajectories, or human videos support action generation and transfer.
  \item \textbf{Reward modeling with foundation models}: VLMs and LVLMs are used to produce feedback, reward, or reward benchmarks.
  \item \textbf{Connection to VLA/world models}: works such as OpenVLA \citep{kim2024openvla}, OpenVLA-OFT \citep{wang2025openvlaoft}, and VLA-RFT \citep{li2025vlarft} show a trend toward integrating reward, RL, and imagined rollouts more closely with VLA systems.
\end{enumerate}

This proposal sits at the intersection of directions (2) and (3),
while preserving a natural upgrade path toward direction (4) once the reward pipeline is stable.
